\documentclass[11pt]{article}
\usepackage[margin=1in]{geometry}
\usepackage{enumitem}
% =============================================================================
% Preambles/header.tex — shared LaTeX/Beamer preamble for this project
%
% Use from a Beamer lecture by prepending:
%     \input{header}
% and compiling with TEXINPUTS=../Preambles:$TEXINPUTS (the /compile-latex
% skill does this automatically).
%
% PALETTE CONTRACT. Color names defined here MUST match the names in
% Quarto/theme-template.scss so Beamer and Quarto renderings use the same
% palette. The scripts/check-palette-sync.sh script enforces this.
%
% To customize: edit the HEX values below AND the matching $variable in
% Quarto/theme-template.scss, then run scripts/check-palette-sync.sh.
% =============================================================================

% -----------------------------------------------------------------------------
% Engine detection + encoding
%
% Our /compile-latex skill uses XeLaTeX, where inputenc is unnecessary and
% can produce encoding warnings. Only load inputenc under pdfTeX.
% -----------------------------------------------------------------------------
\usepackage{iftex}
\ifPDFTeX
  \usepackage[utf8]{inputenc}
\fi

\usepackage{amsmath, amssymb, amsthm}
\usepackage{booktabs}
\usepackage{graphicx}

% -----------------------------------------------------------------------------
% PALETTE — mirrors Quarto/theme-template.scss variable names.
% Load xcolor and define colors BEFORE anything that references them
% (notably hyperref, hypersetup, and the Beamer color assignments).
% -----------------------------------------------------------------------------
\usepackage{xcolor}

\definecolor{primary-blue}{HTML}{012169}      % headings, accents
\definecolor{primary-gold}{HTML}{B9975B}      % emphasis, borders
\definecolor{highlight-yellow}{HTML}{F2A900}  % markers, alerts
\definecolor{light-bg}{HTML}{E8EDF5}          % title / section backgrounds
\definecolor{jet}{HTML}{1A1A1A}               % body text

% Semantic colors (used in diagrams and annotations)
\definecolor{positive}{HTML}{15803D}          % good / observed / identified
\definecolor{negative}{HTML}{B91C1C}          % bad / problematic / bias
\definecolor{neutral}{HTML}{525252}           % reference / context

% Highlight accents (match .hi-* classes in the SCSS)
\definecolor{hi-slate}{HTML}{314F4F}
\definecolor{hi-green}{HTML}{2E8B57}
\definecolor{hi-red}{HTML}{C41E3A}

% Semantic aliases so skill templates and snippets can write
% \textcolor{accent}{...} without hard-coding "primary-blue".
\colorlet{accent}{primary-blue}
\colorlet{accent2}{primary-gold}

% -----------------------------------------------------------------------------
% Hyperref — loaded AFTER xcolor + palette so link colors resolve.
% -----------------------------------------------------------------------------
\usepackage{hyperref}
\hypersetup{colorlinks=true, linkcolor=primary-blue, urlcolor=primary-blue,
            citecolor=primary-blue, pdfborder={0 0 0}}

% -----------------------------------------------------------------------------
% Course code — set once per deck (after \input{header}, before \begin{document})
% via \coursecode{CS401}. Shown in the footer on every slide so a deck's course
% is identifiable without opening the title page. Empty by default.
% -----------------------------------------------------------------------------
\makeatletter
\newcommand{\coursecode}[1]{\gdef\@coursecodetext{#1}}
\gdef\@coursecodetext{}
\makeatother

% -----------------------------------------------------------------------------
% Beamer theme assignments (only apply under Beamer).
%
% We detect Beamer via \@ifclassloaded{beamer}, which is the documented,
% non-internal way. This must be wrapped in \makeatletter/\makeatother
% because \@ifclassloaded contains an @-catcode character.
% -----------------------------------------------------------------------------
\makeatletter
\@ifclassloaded{beamer}{%
  \usefonttheme{professionalfonts}
  \setbeamercolor{structure}{fg=primary-blue}
  \setbeamercolor{frametitle}{fg=primary-blue}
  \setbeamercolor{title}{fg=primary-blue}
  \setbeamercolor{subtitle}{fg=primary-gold}
  \setbeamercolor{author}{fg=primary-blue}
  \setbeamercolor{institute}{fg=neutral}
  \setbeamercolor{date}{fg=primary-gold}
  \setbeamercolor{itemize item}{fg=primary-blue}
  \setbeamercolor{itemize subitem}{fg=primary-gold}
  \setbeamercolor{alerted text}{fg=primary-gold}
  \setbeamercolor{block title}{fg=white, bg=primary-blue}
  \setbeamercolor{block body}{bg=light-bg}
  % Semantic block variants: block=definition/concept (blue), exampleblock=
  % worked example (green/positive), alertblock=key takeaway or caution (gold).
  % Keep to <=2 colored boxes per slide across ALL three variants (INV-7).
  \setbeamercolor{block title example}{fg=white, bg=positive}
  \setbeamercolor{block body example}{bg=light-bg}
  \setbeamercolor{block title alerted}{fg=white, bg=primary-gold}
  \setbeamercolor{block body alerted}{bg=light-bg}

  % Minimal footer: course code (left) + page X / N (right), no navigation clutter
  \setbeamertemplate{navigation symbols}{}
  \setbeamertemplate{footline}{%
    \color{neutral}\footnotesize\hspace{0.7em}\@coursecodetext\hfill
    \insertframenumber/\inserttotalframenumber\hspace{0.7em}\vspace{0.3em}}
}{}
\makeatother

% -----------------------------------------------------------------------------
% TikZ setup — libraries the snippet gallery uses
% -----------------------------------------------------------------------------
\usepackage{tikz}
\usetikzlibrary{arrows.meta, positioning, calc, decorations.pathreplacing,
                fit, shapes.geometric, backgrounds}

% Shared TikZ styles — snippets and hand-written diagrams can pick these up
% via `\begin{tikzpicture}[every node/.style={...}]` or by naming specific
% styles (e.g., `\node[dag-node] (X) at (0,0) {X};`).
\tikzset{
  % Node styles for causal DAGs and similar
  dag-node/.style = {
    draw=primary-blue, fill=light-bg, rounded corners=2pt,
    minimum width=1.2cm, minimum height=0.9cm, align=center, font=\small
  },
  decision-node/.style = {
    draw=primary-gold, fill=white, diamond, aspect=1.8,
    minimum width=1.8cm, minimum height=1.2cm, align=center, font=\small,
    inner sep=1pt
  },
  % Edge styles — observed vs counterfactual
  observed-edge/.style = {-{Latex[length=2.5mm]}, thick, draw=jet},
  counterfactual-edge/.style = {-{Latex[length=2.5mm]}, thick, draw=neutral, dashed},
  confound-edge/.style = {-{Latex[length=2.5mm]}, thick, draw=negative, dashed},
  % Data-point styles for plots
  observed-dot/.style = {fill=primary-blue, draw=primary-blue, circle, inner sep=1.2pt},
  counterfactual-dot/.style = {fill=white, draw=neutral, circle, inner sep=1.2pt}
}

% -----------------------------------------------------------------------------
% Convenience macros
% -----------------------------------------------------------------------------
\newcommand{\muted}[1]{\textcolor{neutral}{#1}}
\newcommand{\key}[1]{\textcolor{primary-gold}{\textbf{#1}}}
\newcommand{\good}[1]{\textcolor{positive}{#1}}
\newcommand{\bad}[1]{\textcolor{negative}{#1}}

% Standout frame macro: full-bleed dark background for section transitions.
% Beamer-only (uses \begin{frame}/\setbeamercolor/\usebeamerfont) -- do not
% call from an article-class file (e.g. Notes/<CODE>/*-notes.tex). \muted,
% \key, \good, \bad, and \coursecode above are plain xcolor/gdef and are
% safe to use from either class; verified when Notes/article-class support
% was added.
% Expand: \transitionslide{Your transition title}
\newcommand{\transitionslide}[1]{%
  \begingroup
  \setbeamercolor{background canvas}{bg=primary-blue}
  \begin{frame}[plain]
    \centering
    \vfill
    {\usebeamerfont{title}\color{white}\Large #1\par}
    \vfill
  \end{frame}
  \endgroup
}

% Section divider frame (Beamer-only), an alternative to \transitionslide with a
% darker two-line style: near-black (jet) background, a small white label line
% above a larger gold title, and a thin gold rule along the bottom edge.
% Expand: \sectiondivider{Section 2}{Pointers}
\newcommand{\sectiondivider}[2]{%
  \begingroup
  \setbeamercolor{background canvas}{bg=jet}
  \begin{frame}[plain]
    \vspace*{3.0cm}
    {\color{white}\ttfamily\large #1\par}
    \vspace{0.5cm}
    {\color{primary-gold}\LARGE #2\par}
    \begin{tikzpicture}[remember picture, overlay]
      \draw[primary-gold, line width=2.5pt]
        ($(current page.south west)+(0,0.1)$) -- ($(current page.south east)+(0,0.1)$);
    \end{tikzpicture}
  \end{frame}
  \endgroup
}

% -----------------------------------------------------------------------------
% Channel watermark — "Concepts That Click" (YouTube).
%
% Article-class files (CompetitiveExam Q/A sets, Notes, Assignments) get a
% light diagonal draft watermark on every page plus a centered footer naming
% the channel. Beamer decks get a subtle TikZ background watermark instead
% (the footline already carries the course code). Centralized here so every
% MCQ PDF is branded without per-file edits.
% -----------------------------------------------------------------------------
\makeatletter
\@ifclassloaded{article}{%
  \usepackage{draftwatermark}
  \SetWatermarkText{Concepts That Click}
  \SetWatermarkScale{1}
  \SetWatermarkFontSize{2cm}
  \SetWatermarkLightness{0.86}
  \SetWatermarkAngle{45}
  \usepackage{fancyhdr}
  \pagestyle{fancy}
  \fancyhf{}
  \fancyfoot[C]{\footnotesize\textcolor{neutral}{Concepts That Click~|~YouTube: \href{https://www.youtube.com/@concepts.thatclick}{@concepts.thatclick}}}
  \renewcommand{\headrulewidth}{0pt}
  \renewcommand{\footrulewidth}{0pt}
  \fancypagestyle{plain}{%
    \fancyhf{}%
    \fancyfoot[C]{\footnotesize\textcolor{neutral}{Concepts That Click~|~YouTube: \href{https://www.youtube.com/@concepts.thatclick}{@concepts.thatclick}}}%
    \renewcommand{\headrulewidth}{0pt}%
    \renewcommand{\footrulewidth}{0pt}%
  }
}{}
\@ifclassloaded{beamer}{%
  \setbeamertemplate{background}{%
    \begin{tikzpicture}[remember picture, overlay]
      \node[rotate=45, opacity=0.055, align=center]
        at (current page.center)
        {\Huge\textcolor{neutral}{Concepts That Click}};
    \end{tikzpicture}%
  }
}{}
\makeatother 
\coursecode{JKSSB}
\renewcommand{\thesection}{65.\arabic{section}}
\newtheorem{example}{Example}[section]
\newenvironment{definitionbox}[1]{%
  \par\noindent\textbf{Definition (#1).}\ \itshape
}{\par\vspace{0.3em}}
\title{Current Multimedia, GUI and Media Formats --- Detailed Lecture Notes}
\author{JKSSB Computer Awareness}
\date{\today}

\begin{document}
\maketitle

\section{What multimedia means}
\textbf{Multimedia} is the presentation of content that combines more than
one medium in a single document or program. The media may be \textbf{text,
graphics, audio, video}, and \textbf{animation}. Text and graphics are
\emph{static} media: they do not change with time. Audio, video, and animation
are \emph{time-based} media: their meaning depends on being played in the
correct order and at the correct speed. Because a computer stores everything
as binary data, every medium --- sound, picture, or moving image --- must
first be captured and converted into numbers before it can be stored,
compressed, and replayed.

\begin{definitionbox}{Multimedia}
Content that combines two or more media types --- such as text, graphics,
audio, video, and animation --- in a single digital presentation.
\end{definitionbox}

\begin{center}
\begin{tabular}{p{0.20\textwidth}p{0.68\textwidth}}
\toprule
\textbf{Element} & \textbf{Meaning} \\
\midrule
Text & Readable characters; the base medium of almost every document. \\
Graphics & Still images, drawings, and photographs. \\
Audio & Sound, music, and speech. \\
Video & Moving pictures, usually accompanied by sound. \\
Animation & A sequence of images played fast enough to create motion. \\
\bottomrule
\end{tabular}
\end{center}

\section{Digitising media: sampling and quantization}
An analogue signal, such as a sound wave or a scene's light, is
\emph{continuous}: it has a value at every instant. A computer, however,
stores only discrete numbers. Converting a continuous signal into digital
form is done by \textbf{analogue-to-digital conversion}, and it has two
distinct steps, \textbf{sampling} and \textbf{quantization}.

\begin{definitionbox}{Sampling}
The process of measuring a continuous signal at regular intervals and
recording each measurement. The number of measurements taken per second is
the \textbf{sampling rate}, measured in hertz (Hz).
\end{definitionbox}

A \textbf{higher sampling rate} captures the signal more often, so the
reproduction is more faithful, but the file becomes larger. A lower sampling
rate saves space but loses detail. The sampling rate is the horizontal
resolution of the digitised signal: how often we look at it.

\begin{definitionbox}{Quantization}
The process of mapping each measured sample to one of a fixed set of discrete
levels, because a computer can store only a limited number of distinct values
per sample.
\end{definitionbox}

The number of available levels is set by the \textbf{bit depth} of each
sample. With a higher bit depth --- more bits per sample --- there are more
levels, so each recorded value is closer to the true value; the rounding
error, called \emph{quantization noise}, is smaller, and quality is better,
but again the file is larger.

\begin{example}
Think of drawing a smooth curve on graph paper. \textbf{Sampling} decides how
many points across the page you mark --- more points follow the curve more
closely. \textbf{Quantization} decides how finely the vertical axis is ruled
--- more lines let you record each point's height more precisely. The two
choices together determine how closely the digitised version matches the
original, and how big the file becomes.
\end{example}

The common exam confusion is to swap the two. Remember the one-line
distinction: \textbf{sampling is how often we measure; quantization is how
precisely each measurement is stored.}

\section{Compression: lossy and lossless}
Raw multimedia files are large, so they are almost always \textbf{compressed}
--- stored in a form that takes less space. There are two kinds of
compression, and the difference is whether the original can be recovered
exactly.

\begin{definitionbox}{Lossless compression}
Compression that allows the original data to be reconstructed exactly. No
information is discarded. PNG and GIF images and ZIP archives are lossless.
\end{definitionbox}

\begin{definitionbox}{Lossy compression}
Compression that discards some data to reduce the file size further, so the
reconstruction is only an approximation of the original. JPEG images and MP3
audio are lossy.
\end{definitionbox}

Lossless compression works by removing redundancy --- repeated or predictable
data --- while keeping every value. Lossy compression goes further: it removes
detail that a human eye or ear is unlikely to notice. Lossy files are usually
much smaller, which is why lossy formats dominate the internet for
photographs, music, and video. When exactness matters --- for a diagram,
screenshot, or archive --- a lossless format is used instead.

\section{Synchronization}
A multimedia file often carries more than one \textbf{stream} at once: a video
stream and an audio stream, for instance. \textbf{Synchronization} is the
process of keeping those streams aligned in time so they play together
correctly. If the streams drift out of step, the sound no longer matches the
picture --- the narrator's voice is heard before or after the lips move, which
is called a loss of \emph{sync}. Synchronization is what makes a video play
with the sound in the right place at the right moment; it is a timing
requirement peculiar to time-based media.

\begin{definitionbox}{Synchronization}
The alignment in time of separate media streams --- such as audio and video
--- so that they are presented together correctly during playback.
\end{definitionbox}

\section{The graphical user interface and its elements}
A \textbf{GUI} is a \textbf{Graphical User Interface}: it lets the user work
with pictures, icons, and menus instead of typing commands. A \textbf{CLI} is
a \textbf{Command-Line Interface}: it accepts typed text commands. The GUI is
easier for beginners and is the standard on modern desktops; the CLI is often
faster for an expert and is still used for administration.

\begin{definitionbox}{GUI (Graphical User Interface)}
A user interface in which the user interacts with the computer through visual
elements such as windows, icons, menus, and a pointer, rather than by typing
text commands.
\end{definitionbox}

The four classic GUI elements are remembered by the short form \textbf{WIMP}:
\textbf{Windows, Icons, Menus, Pointers}.

\begin{center}
\begin{tabular}{p{0.16\textwidth}p{0.72\textwidth}}
\toprule
\textbf{Element} & \textbf{Meaning} \\
\midrule
Window & A rectangular on-screen area that holds a program or a document. \\
Icon & A small picture that represents a file, folder, or program. \\
Menu & A list of commands that the software offers. \\
Pointer & The on-screen arrow, moved by the mouse or touchpad, used to select. \\
\bottomrule
\end{tabular}
\end{center}

\begin{definitionbox}{Dialogue box}
A small window that appears to ask the user for input or to confirm an action,
such as Save As or Print. It usually contains text boxes, drop-down lists,
check boxes, radio buttons, and OK/Cancel buttons.
\end{definitionbox}

\section{WYSIWYG}
\textbf{WYSIWYG} stands for \textbf{What You See Is What You Get}. It means
that what is shown on the screen is exactly how the document will look when
printed --- the same layout, fonts, spacing, and page breaks. A word processor
such as MS Word works in WYSIWYG mode, so the page visible on the monitor
matches the printed page.

\begin{example}
While editing a document in MS Word, the text, margins, and page breaks on
screen match the printed output. This is WYSIWYG. When the user opens the
Print dialogue box and chooses options, those settings are entered in a
dialogue box, not typed on the command line. The two ideas are often asked
together: WYSIWYG describes the on-screen view; a dialogue box is one GUI
element that collects input.
\end{example}

\section{Image formats: PNG, JPEG, and GIF}
Three image formats appear most often in competitive-exam questions.

\begin{center}
\begin{tabular}{p{0.12\textwidth}p{0.26\textwidth}p{0.50\textwidth}}
\toprule
\textbf{Format} & \textbf{Full form} & \textbf{Nature and typical use} \\
\midrule
PNG & Portable Network Graphics & Lossless; supports transparency; used for logos and web graphics. \\
JPEG & Joint Photographic Experts Group & Lossy; best for photographs and complex pictures. \\
GIF & Graphics Interchange Format & Lossless but limited to 256 colours; used for simple graphics and short animation. \\
\bottomrule
\end{tabular}
\end{center}

\textbf{JPEG} is chosen for photographs because its lossy compression makes
them much smaller with little visible loss. \textbf{PNG} is chosen when exact
colours or transparency are needed, because it is lossless. \textbf{GIF} is
chosen for simple graphics and small animations, but it can show only 256
colours, so it is unsuitable for photographs. \textbf{BMP} (Bitmap) is an
older, largely uncompressed image format that produces large files and is not
used for the web.

\section{Audio and video formats: MP3 and MP4}
\textbf{MP3} stands for \textbf{MPEG-1 Audio Layer III} (also written MPEG
Audio Layer III). It is a \textbf{lossy} audio format: it removes sound detail
that the human ear is unlikely to notice, which is why an MP3 file is far
smaller than an uncompressed one. Its small size made MP3 the standard for
portable music and downloads. By contrast, \textbf{WAV} is an uncompressed
audio format.

\textbf{MP4} stands for \textbf{MPEG-4 Part 14}. It is a \textbf{container}
format: a single MP4 file can hold video, audio, and subtitles together. It is
widely used for streaming, downloads, and mobile video. A container such as
MP4 holds the streams; the \emph{codec} inside the container is what actually
compresses the video and audio. \textbf{AVI} (Audio Video Interleave) is an
older video container.

\begin{definitionbox}{MPEG}
MPEG stands for Moving Picture Experts Group --- the group that defines the
MP3 audio layer and the MP4 container, among other standards.
\end{definitionbox}

\section{Frame rate, resolution, and bit rate}
Video is a rapid sequence of still frames. The \textbf{frame rate} is the
number of frames shown per second (fps); a higher frame rate produces smoother
motion but a larger file. The \textbf{resolution} is the number of pixels in
each frame (for example, $1920 \times 1080$); more pixels mean a sharper image
and a larger file. The \textbf{bit rate} is the amount of data stored or sent
per second; a higher bit rate keeps more detail but needs more space or
bandwidth. These three settings, together with the choice of lossy or lossless
compression, decide the final quality and size of a video or audio file.

\begin{definitionbox}{Frame rate}
The number of still frames displayed per second in a video. A higher frame
rate gives smoother motion at the cost of a larger file.
\end{definitionbox}

\begin{example}
A video at 30 frames per second with a high bit rate looks smooth and sharp
but takes much more space than the same clip at a lower frame rate and bit
rate. Streaming services lower the bit rate when the network is slow, trading
quality for smooth playback --- a daily example of the quality-versus-size
trade-off that runs through the whole lesson.
\end{example}

\section{Plotter}
A \textbf{plotter} is an \textbf{output} device that draws on paper. Unlike a
printer, which lays down text and small images, a plotter produces
\textbf{vector graphics}: lines and curves defined by coordinates rather than
by a grid of dots. It is used for \textbf{large engineering drawings}, maps,
blueprints, and banners, where continuous, accurate lines are needed. Early
pen plotters moved a pen across the paper; modern inkjet plotters spray ink
while the sheet moves. A plotter is always an output device, never an input
device (PYQ-29, PYQ-08).

\section{Impact versus non-impact printers (revisit)}
The basic classification of printers reappears in later units, so keep it
firmly in mind. An \textbf{impact printer} forms an image by striking the
paper; a \textbf{non-impact printer} does so without any striking action.

\begin{center}
\begin{tabular}{p{0.16\textwidth}p{0.36\textwidth}p{0.36\textwidth}}
\toprule
\textbf{Point} & \textbf{Impact} & \textbf{Non-impact} \\
\midrule
Action & Strikes the paper to print & Forms the image without striking \\
Examples & Dot-matrix, daisy-wheel, line printer & Inkjet, laser, thermal printer \\
Carbon copies & Possible (same strike hits several sheets) & Not possible \\
Noise & Comparatively noisy & Comparatively quiet \\
\bottomrule
\end{tabular}
\end{center}

The quick test is the mechanism: if the device physically hits the paper, it
is impact; if it does not, it is non-impact. Both kinds are output devices. A
\emph{dot-matrix} printer is impact; an \emph{inkjet} and a \emph{laser}
printer are non-impact.

\section{Multimedia hardware and applications}
Playing and creating multimedia needs both hardware and software. On the
hardware side, a \textbf{sound card} processes audio, \textbf{speakers} and
\textbf{headphones} play it, and a \textbf{microphone} captures it. A
\textbf{graphics card (GPU)} renders images and video, a \textbf{monitor or
projector} displays them, and a \textbf{camera} or \textbf{webcam} captures
still and moving pictures. Because audio and video files are large, multimedia
also places heavy demands on \textbf{memory and storage}.

On the software side, editors for images, audio, and video are used to
capture, edit, compress, and synchronize the media streams. The main
application areas of multimedia are education (e-learning and tutorials),
entertainment (music, games, and streaming films), business (presentations
and advertising), and medicine (medical imaging and simulation).

\begin{center}
\begin{tabular}{p{0.22\textwidth}p{0.66\textwidth}}
\toprule
\textbf{Short form} & \textbf{Full form} \\
\midrule
GUI & Graphical User Interface \\
CLI & Command-Line Interface \\
WYSIWYG & What You See Is What You Get \\
WIMP & Windows, Icons, Menus, Pointers \\
PNG & Portable Network Graphics \\
JPEG & Joint Photographic Experts Group \\
GIF & Graphics Interchange Format \\
MP3 & MPEG-1 Audio Layer III \\
MP4 & MPEG-4 Part 14 \\
MPEG & Moving Picture Experts Group \\
AVI & Audio Video Interleave \\
BMP & Bitmap \\
\bottomrule
\end{tabular}
\end{center}

\section{Exam retrieval and common confusions}
Several distinctions prevent most errors on this topic.
\begin{enumerate}
  \item \textbf{Sampling} is how often a signal is measured; \textbf{quantization}
    is how precisely each sample is stored. Do not swap them.
  \item \textbf{Lossless} formats restore the original exactly (PNG, GIF);
    \textbf{lossy} formats discard data (JPEG, MP3, MP4).
  \item \textbf{JPEG} is lossy and best for photographs; \textbf{PNG} is
    lossless and supports transparency; \textbf{GIF} is lossless and supports
    simple animation.
  \item \textbf{MP3} is audio (MPEG-1 Audio Layer III); \textbf{MP4} is video
    (MPEG-4 Part 14). Do not swap the full forms.
  \item \textbf{Synchronization} keeps audio and video aligned in time; it is
    not compression and not sampling.
  \item \textbf{GUI} is Graphical User Interface; its elements are
    \textbf{Windows, Icons, Menus, Pointers} (WIMP).
  \item \textbf{WYSIWYG} means the screen matches the print; a \textbf{dialogue
    box} is a small window that collects input.
  \item A \textbf{plotter} is an output device for large drawings; an
    \textbf{impact} printer strikes the paper and a \textbf{non-impact} printer
    does not.
\end{enumerate}

\begin{example}
An item asks, ``Which format is lossless and supports transparency?'' Trace
the formats: JPEG is lossy, MP3 is lossy audio, MP4 is lossy video, and GIF is
lossless but limited to 256 colours without transparency support. The answer
is PNG, the lossless format that supports transparency.
\end{example}

\subsection*{Exercises}
\begin{enumerate}[label=\arabic*.]
  \item Define multimedia and name its five building blocks.
  \item Distinguish sampling from quantization, and explain how each affects
    file size.
  \item Differentiate lossy and lossless compression, giving one example of
    each.
  \item What is synchronization in multimedia, and why is it needed?
  \item Expand GUI, WYSIWYG, and WIMP.
  \item State one suitable use each for PNG, JPEG, and GIF.
  \item Expand MP3 and MP4, and state the medium each is used for.
  \item Explain why a plotter is preferred over a printer for large engineering
    drawings.
  \item Distinguish impact printers from non-impact printers with one example
    of each.
\end{enumerate}

\subsection*{Solutions}
\begingroup\small
\begin{enumerate}[label=\arabic*.]
  \item Multimedia is content combining more than one medium in a single
    presentation. Its five building blocks are text, graphics, audio, video,
    and animation.
  \item Sampling measures a continuous signal at regular intervals; a higher
    sampling rate captures more detail but enlarges the file. Quantization maps
    each sample to one of a fixed set of levels set by the bit depth; more
    levels improve quality but also enlarge the file.
  \item Lossless compression restores the original exactly (for example PNG);
    lossy compression discards some data and is only an approximation (for
    example JPEG).
  \item Synchronization is the alignment in time of separate media streams such
    as audio and video. It is needed so that sound and picture play together
    correctly rather than drifting apart.
  \item GUI = Graphical User Interface; WYSIWYG = What You See Is What You Get;
    WIMP = Windows, Icons, Menus, Pointers.
  \item PNG: logos and web graphics needing transparency or exact colours.
    JPEG: photographs and complex pictures. GIF: simple graphics and short
    animations.
  \item MP3 = MPEG-1 Audio Layer III, used for audio (music and speech). MP4 =
    MPEG-4 Part 14, used for video (a container holding video, audio, and
    subtitles).
  \item A plotter produces vector graphics --- continuous lines defined by
    coordinates --- on large sheets, which suits precise engineering drawings,
    maps, and blueprints better than a printer.
  \item Impact printers strike the paper (for example dot-matrix); non-impact
    printers form the image without striking (for example inkjet or laser).
\end{enumerate}
\endgroup
\end{document}
